\textbf{Lombard \& Piraux 2004 --- Eq 43--44 (time-stepping at irregular points, paper p.\ 16)}

\medskip

\textbf{Setup.} Once all modified values $\boldsymbol{U}^*_{I,J}$
have been computed at $t^n$ via Eq 42, the irregular points are
updated using the \emph{same} discrete operator $\boldsymbol{H}$
as the regular points (Eq 12 elsewhere in the paper) but with
substituted neighbour values: numerical $\boldsymbol{U}^n$ values
from same-side neighbours, modified $\boldsymbol{U}^*$ values from
across-$\Gamma$ neighbours.

\medskip

\textbf{Eq 43 --- per-stencil-leg substitution rule.}
Denote $\Omega(i, j)$ the medium to which the grid point $(i, j)$
belongs. For each $(\alpha, \beta)$ offset in the stencil:
\begin{equation*}
\begin{aligned}
  \Omega(i + \alpha, j + \beta) = \Omega(i, j)
    &\quad \Rightarrow \quad
    \tilde{\boldsymbol{U}}_{i+\alpha, j+\beta} = \boldsymbol{U}^n_{i+\alpha, j+\beta} \\
  \Omega(i + \alpha, j + \beta) \neq \Omega(i, j)
    &\quad \Rightarrow \quad
    \tilde{\boldsymbol{U}}_{i+\alpha, j+\beta} = \boldsymbol{U}^*_{i+\alpha, j+\beta}
\end{aligned}
\tag{43}
\end{equation*}

\textbf{Eq 44 --- modified update for an irregular point $(i, j)$:}
\begin{equation*}
  \boxed{\ \boldsymbol{U}^{n+1}_{i, j}
  = \boldsymbol{H} \!\left( \tilde{\boldsymbol{U}}_{i+\alpha, j+\beta} \right)\ }
  \tag{44}
\end{equation*}

\medskip

\textbf{Key properties (paper's own framing):}
\begin{itemize}
  \item \emph{Scheme-agnostic.} Coupling the ESIM with ``a wide
    class of scheme is automatic: no modification of the scheme
    is required, and the computation of $\boldsymbol{U}^*_{i,j}$
    does not depend on the discrete operator $\boldsymbol{H}$.''
  \item \emph{Sub-cell resolution.} The ESIM encodes the geometry
    of $\Gamma$: the interface conditions (Eq 19) involve $x'$,
    $y'$ and their successive derivatives up to $k$-th order at
    $P$. The Taylor expansions (Eq 35) and (Eq 37) introduce a
    sub-cell resolution for the position of $P$ within the mesh.
\end{itemize}

\medskip

\textbf{Implementation pattern for Devito (dipping-line Petrobras):}
\begin{enumerate}
  \item \textbf{Identify irregular points} (set-up phase, once):
    grid points $(i, j)$ whose RSG diagonal stencil
    $(i \pm 1, j \pm 1)$ crosses $\Gamma$. Already implemented
    in \texttt{scripts/it\_e\_tier1\_geometry.py}.
  \item \textbf{Precompute the master matrix} per irregular point:
    $\boldsymbol{N}_{I, J} := \boldsymbol{\Pi}_{I, J}^k
    \boldsymbol{G}_1^k \boldsymbol{K}_1^k \bar{\boldsymbol{M}}^{-1}$
    (one fixed row vector). Because $n$, $t$, $\alpha_1$,
    $\alpha_2$ are constant along the dipping line, the four
    sub-matrices are global; $\bar{\boldsymbol{M}}^{-1}$ varies
    only by the integer cell offset of $B$ relative to $P$.
  \item \textbf{Each timestep:}
    \begin{enumerate}
      \item Gather $(\boldsymbol{U}^n)_B$ from the disk
        of radius $q$ around each $P$ on the regular grid.
      \item Compute $\boldsymbol{U}^*_{I, J}
        = \boldsymbol{N}_{I, J} \cdot (\boldsymbol{U}^n)_B$.
      \item Apply $\boldsymbol{H}$ at each irregular point using
        the $\tilde{\boldsymbol{U}}$ substitution rule (Eq 43).
    \end{enumerate}
\end{enumerate}

\medskip

\textbf{Note on sides.} For an irregular point on the $\Omega_1$
side, the master matrix has $\boldsymbol{G}_2^k \boldsymbol{K}_2^k$
instead of $\boldsymbol{G}_1^k \boldsymbol{K}_1^k$ (Eq 42's swap).
Both sides have their own irregular-point list and their own
precomputed $\boldsymbol{N}$ row vectors.
